\begin{textsample}{NEG Dim 9 – global\_north\_2024\_02 – Score -6.00 – global\_north\_2024\_02/106113140.txt}  \label{ex:f9_neg_286}
DUP MP Ian Paisley has been criticised for his heated contribution to the Gaza ceasefire debate in Westminster, prompting a warning from the Deputy Speaker.
The House of Commons later descended into chaos as the Government and SNP condemned Speaker Sir Lindsay Hoyle for his handling of the debate.
Labour ' s amendment pushing for an immediate Gaza ceasefire was approved by the Commons.
Earlier, SNP MPs and some Conservatives walked out of the chamber in an apparent protest at the state of affairs as the debate reached its conclusion.
Commons Leader Penny Mordaunt claimed Sir Lindsay had"hijacked"the debate and"undermined the confidence"of the House in its longstanding rules by selecting Labour ' s bid to amend the SNP motion calling for an"immediate ceasefire"in Gaza and Israel.
It had been expected that Sir Lindsay would select just the Government ' s amendment seeking an"immediate humanitarian pause"to the Israel-Hamas conflict, which could pave the way for a more permanent stop in fighting.
But instead, he decided that the Commons would first @ @ @ @ @ @ @ @ @ @ ceasefire"before moving on to further votes on the SNP ' s original motion, and then the Government ' s proposals if either of the first two were to fail to garner enough support.
Conservative MP William Wragg, who called for Sir Lindsay to resign, later tried to make the House of Commons sit in private.
SNP MPs were understood to have headed to the voting lobby after the walkout from the chamber.
The debate sparked passions on all sides.
DUP MP Mr Paisley recalled a recent meeting with Adi Efrat, who was \textbf{kidnapped} by Hamas on October 7 after her house was scorched to the ground.
The 51-year-old was taken from \textbf{Kibbutz} Be'eri where 120 neighbours -- including men, women and children -- were killed during the murderous rampage.
Mr Paisley branded the motion -- which he said doesn't contain a single word"about the rape of women and \textbf{murder} of children"-- as"vile"."It ' s as if, Mr Deputy Speaker, it did @ @ @ @ @ @ @ @ @ @ ' s as if it was invisible."It ' s as if like other people in the 20th century who denied what happened."It is utterly vile."The North Antrim representative said the"awful"war in Gaza is a consequence of"the unjustifiable attack"on Israelis and Jews on October 7.
SDLP MP Claire Hanna, who previously worked for a relief and development agency, told the Commons she has never encountered a humanitarian crisis as"hellish"as what the people of Gaza are facing."Our constituents are watching in devastation and distress and feel powerless to the point of complicity,"she said.
While stressing that she stood in"full solidarity"with the Israeli victims of the Hamas attack, the South Belfast MP said the scenes coming from Gaza filled people with dread."I feel that way,"she continued."I look at my six-year-old sweet, smiling innocent \textbf{daughter}, and I see a six year @ @ @ @ @ @ @ @ @ @ her cries and the bodies of all the people she \textbf{loved} -- those scenes will never leave people."\textbf{Ms} Hanna objected to the"slurs and distortions"put on those who show empathy as she branded the October 7 massacre as"wicked"and"vile".
She described Hamas as a"cynical organisation"and said she believes in Israel ' s right to exist."Comparisons between the Middle East and NI are shallow,"she added."But the one lesson that can be learned is the first step is stopping the killing... if we don't support a ceasefire now, when will we?"Later, Sir Lindsay apologised to the Commons amid shouts of"resign"from some MPs.
Sir Lindsay told the Commons :"I thought I was doing the right thing and the best thing, and I regret it, and I apologise for how it ' s ended up."I do take responsibility for my actions, and that @ @ @ @ @ @ @ @ @ @ who have been involved."

% matched lemmas: daughter, kibbutz, kidnap, love, ms, murder
\end{textsample}
